\section{Mục tiêu}
Mục tiêu trọng tâm của nghiên cứu này là đề xuất và triển khai 
một phương pháp tiếp cận mới dựa trên Đồ thị các Tập Lồi (Graphs of Convex Sets - GCS) 
nhằm giải quyết bài toán hoạch định chuyển động cho robot trong môi trường chứa vật cản. 
Nghiên cứu hướng đến việc khắc phục những hạn chế về tính tối ưu và khả năng xử lý các ràng
buộc động lực học phức tạp thường gặp ở các phương pháp lấy mẫu truyền thống bằng cách phân hoạch 
không gian tự do thành các vùng lồi hữu hạn và tham số hóa quỹ đạo sử dụng đường cong Bézier. 
Luận văn đặt mục tiêu thiết lập bài toán dưới dạng tối ưu hóa hỗn hợp nguyên, 
sau đó áp dụng kỹ thuật nới lỏng lồi (convex relaxation) chặt chẽ để chuyển đổi thành bài toán lồi chuẩn,
nhằm đảm bảo tìm được quỹ đạo tối ưu toàn cục với thời gian tính toán hiệu quả. 
Cuối cùng, nghiên cứu tập trung chứng minh tính vượt trội của phương pháp GCS so với các thuật toán dựa trên phương pháp lấy mẫu như PRM hay RRT 
thông qua việc giảm thiểu thời gian truy vấn và rút ngắn độ dài quỹ đạo trên các hệ thống có số chiều cao.